\subsection{Pressure-induced shear behind an isotropic mineral}
\label{sec:fabric-pressure-shear}

A pressure increment in a specimen held at fixed deformation requires
boundary reactions. Their shear component provides a direct test of the
nonspherical Biot tensor. We use an isotropic mineral with
\(\mathcal{K}_s=2.5K_*\), shear modulus \(5K_*/6\), and
\(\phi_{s0}=0.9\). The retained distention block of
\cref{sec:fabric-biot} has \(k_v=5.4K_*\), \(k_a=K_*\), and
\(k_c=0.4K_*\). This positive definite block couples volume and shape.
The comparison with \(k_c=0\) retains shape compliance; the spherical
comparison freezes shape and retains the same volumetric modulus.

Write the prescribed in-plane axis as
\(\mathbf m=(\cos\theta,\sin\theta,0)^T\). At the reference state,
\eqref{eq:fabric-transverse-biot} and the pressure tangent give
\begin{equation}
 \left.\frac{\partial\sigma_{12}}{\partial p}\right|_{\mathbf F=\mathbf I,p=0}
 =-(B_\parallel-B_\perp)\sin\theta\cos\theta.
 \label{eq:fabric-pressure-shear-reaction}
\end{equation}
Thus the shear reaction reverses with the fabric angle and vanishes when
the axis lies along either coordinate direction. Its largest magnitude
occurs at \(45^\circ\). \Cref{fig:fe-fabric-probe} evaluates the
reference law, with the pressure reactions checked independently by
eliminating the two distention strains. An uncoupled shape compliance can
change the drained stiffness while leaving the pressure reaction
spherical. Volume--shape coupling is what makes this experiment distinguish
the fabric model from the spherical specialization. The modulus sweep
in panel (b) keeps \(k_v\) fixed and varies \(k_a/k_v\) and
\(k_c/\sqrt{k_vk_a}\). The latter lies between \(-0.9\) and
\(0.9\), so every distention block remains positive definite.
The zero reaction at \(k_c=0\) persists across the entire range of
shape stiffness, while its magnitude and sign change with volume--shape
coupling.

\begin{figure}[htbp]
 \centering
 \includegraphics[width=\textwidth]{build/insight/fabric_pressure_shear.pdf}
 \caption{Pressure-induced shear reaction at fixed deformation for an
 isotropic mineral. (a) Reaction per pressure increment against fabric
 orientation. (b) Reaction per pressure increment at \(45^\circ\) as
 shape stiffness and volume--shape coupling vary within the positive
 definite domain. The circle marks the moduli used in panel (a). The
 uncoupled and spherical curves in panel (a) coincide at zero shear reaction. All moduli and states are synthetic.}
 \label{fig:fe-fabric-probe}
\end{figure}

\subsection{Consolidation with tensor and scalar pressure coupling}
\label{sec:fabric-scalar-comparison}

A scalar coefficient calibrated to a mean pressure reaction need not
reproduce a constrained specimen's settlement and lateral reaction.
We compare the reference fabric law with a scalar approximation calibrated
to preserve \(\tr\mathbf B_0/3\). Both use the same drained stiffness,
fixed-strain storage, and isotropic mobility. Solid storage is evaluated
from the reference fabric equilibrium using
\eqref{eq:reference-solid-storage-definition}. The scalar coefficient
replaces the tensor in both stress and fluid mass:
\begin{align}
 \mathbf\sigma_{\mathrm{scalar}}
 &=\mathbb C^d:\mathbf\varepsilon
       -\frac{\tr\mathbf B_0}{3}p\mathbf I,
 \label{eq:fabric-scalar-stress}\\
 \frac{m_{f,\mathrm{scalar}}}{\bar\rho_{f0}}
 &=1-\phi_{s0}+\frac{\tr\mathbf B_0}{3}\tr\mathbf\varepsilon
       +\left(\frac{1-\phi_{s0}}{K_f}+S_s\right)p.
 \label{eq:fabric-scalar-mass}
\end{align}
Here the reference strain is
\(\mathbf\varepsilon=\operatorname{sym}\operatorname{Grad}\mathbf u\).
The tensor model replaces the scalar strain contraction by
\(\mathbf B_0:\mathbf\varepsilon\) and the pressure term in stress by
\(-p\mathbf B_0\). Using the same coupling in both equations preserves
reciprocity. A prescribed-deformation material test checks both compiled
stress and mass against a separate linear solve after this substitution.
For an anisotropic mineral, pressure differentiation of the compiled
mineral volume also agrees with the solid part of the fluid-mass storage.
With positive storage and positive definite drained stiffness,
each model has a positive definite energy in strain and fluid content.
The scalar model is a phenomenological comparison; its internal fabric
variables are not interpreted as mineral-volume predictions.

The plane-strain strip occupies \(0\le X_1\le1\),
\(0\le X_2\le0.1\). Both vertical sides have \(u_1=0\), the bottom
has \(u_2=0\), and the upper surface carries a uniform prescribed vertical traction.
Settlement is the mean vertical displacement of that surface. Tangential
tractions on the horizontal boundaries are zero. Only the right side
drains, with \(p=0\); the remaining boundaries are impermeable.
Starting from zero displacement and pressure, the compressive traction
increases linearly to \(10^{-4}K_*\) at \(t=0.01\) and remains constant
until \(t=0.75\). We use \(K_f=8K_*\), reference fluid density one,
and constant isotropic mobility \(k/\mu_f=1.5\) in the consistent
nondimensional units of the calculation. Lengths are normalized by the
strip width, and the time unit is chosen to give this nondimensional
mobility. The prescribed fabric angles
are \(0^\circ\), \(45^\circ\), and \(90^\circ\); each scalar/tensor
pair has identical mechanical and hydraulic boundary conditions.

\input{build/insight/comparison_summary.tex}
For these parameters and constraints, pressure is more sensitive to the
scalar approximation than settlement. The shared drained stiffness gives
each scalar/tensor pair the same fully drained settlement.

\begin{figure}[htbp]
 \centering
 \includegraphics[width=\textwidth]{build/insight/fabric_comparison.pdf}
 \caption{Consolidation from rest through a load ramp and subsequent
 drainage. Solid curves use the tensor coupling and dashed curves its
 mean scalar approximation, with the drained stiffness, storage and
 mobility held fixed within each orientation. The panels show pressure
 at the sealed lower-left corner, mean settlement of the loaded upper surface, and
 horizontal reaction on the constrained right side. The final panel
 quantifies scalar-approximation errors relative to the tensor histories;
 the norm and refinement comparison are stated in the text.}
 \label{fig:fe-fabric-mandel}
\end{figure}

\subsection{Finite pure shear with equilibrated pore volume and shape}
\label{sec:fabric-finite-benchmark}

To separate finite mineral-volume effects from transport, we prescribe a
homogeneous isochoric pure stretch with principal axes oblique to the
coordinate axes. Let \(\mathbf m\) be the fabric axis at
\(45^\circ\), and let \(\mathbf n\) be its in-plane unit normal.
We prescribe
\begin{equation}
 \mathbf F=\exp\left[\gamma
       (\mathbf m\otimes\mathbf m-\mathbf n\otimes\mathbf n)\right],
 \qquad J=1,
 \qquad -0.6\le\gamma\le0.6.
 \label{eq:fabric-finite-pure-shear}
\end{equation}
The principal stretches are \(e^\gamma,e^{-\gamma},1\). For this
isotropic mineral, the distention, mineral stretch, and mineral stress
commute, satisfying \eqref{eq:fabric-commuting-work}. The off-diagonal
spatial stress is nevertheless nonzero because the common axes are
oblique. This is a finite pure-shear material experiment, separate from
the reference-state finite-element calculations.

We use the same distention moduli as in
\cref{sec:fabric-pressure-shear} and retain the quadratic energy in
\(\mathbf E_{\mathrm{dis}}=\tfrac12\log\mathbf G\), but solve
\eqref{eq:fabric-equilibrium} with the exact exponential mineral volume
\(\bar J=J/a\). Both \(\ln a\) and \(\ln h\) equilibrate at every
state. The comparison evaluates the reference linear constitutive law
on the same logarithmic strain \(\mathbf\varepsilon=\log\mathbf F\).
This convention isolates constitutive truncation from an additional
replacement of logarithmic strain by engineering strain.

\Cref{fig:fabric-finite} compares the responses at \(p=0.5K_*\) and
maps the stress discrepancy over \(0\le p/K_*\le1\). Subtracting the
zero-pressure stress at each fixed deformation isolates the pressure-induced
shear reaction. Although total shear stress is dominated by the elastic
shape response, its pressure-induced part changes with deformation.
Pore volume also differs from the linear prediction even though the total
volume remains fixed. The shape panel reports the pressure-induced change
in \(\ln h\), exposing the volume--shape coupling in the equilibrium.

Independent central differences of the unreduced spectral mineral energy
at the converged states check all nine spatial stress components, with a
maximum Frobenius-norm discrepancy of \(1.1\times10^{-8}K_*\) at the
smallest perturbation. Re-solving the equilibrium under pressure increments
checks the complete pressure tangent to \(1.3\times10^{-11}\).
An independent minimization of that spectral energy reproduces the two
internal strains to \(1.1\times10^{-8}\). Reducing deformation and
pressure together gives stress errors of second order, with measured
orders \(1.9989\)--\(1.9997\), recovering the reference linearization.
All plotted states have positive phase volumes and positive internal
Hessians. These checks establish local internal equilibrium and work
consistency on the commuting family; they do not establish stability
against every deformation mode or a general noncoaxial finite law.

\begin{figure}[htbp]
 \centering
 \includegraphics[width=\textwidth]{build/insight/fabric_finite.pdf}
 \caption{Finite pure shear with oblique fabric, using exact exponential
 volumes and equilibrated volume and shape strains. Panels (a)--(c) use
 \(p=0.5K_*\): pressure-induced shear stress, pore volume per reference
 mixture volume, and pressure-induced shape change. Each pressure-induced
 quantity subtracts the zero-pressure state at the same deformation.
 Panel (d) shows the full stress error of the reference linearization,
 normalized by the mineral bulk modulus. The finite calculation is
 restricted to commuting mineral stress and distention.}
 \label{fig:fabric-finite}
\end{figure}
